% Partie : counterfactuals
% Chapitre : minimal-change-semantics
% Section : antecedent-strengthening

\documentclass[../../../include/open-logic-section]{subfiles}

\begin{document}

\olfileid[fr]{cnt}{min}{agg}

\olsection{Renforcement de l'antécédent}

Le « renforcement de l'antécédent » désigne l'inférence
$!A \lif !C \Entails (!A \land !B) \lif !C$. Elle est valide pour le
conditionnel matériel, mais pas pour les conditionnels
contrefactuels. Supposons qu'il soit vrai que, si je frottais cette
allumette, elle s'enflammerait : ni l'allumette ni la surface du
carnet d'allumettes ne présentent de défaut, je ne casserais pas
l'allumette, etc. Il n'est pourtant pas vrai que, si je frottais
cette allumette dans le vide spatial, elle s'enflammerait. L'inférence
suivante n'est donc pas valide :
\begin{quote}
  Si l'on frottait cette allumette, elle s'enflammerait.

  Donc, si l'on frottait cette allumette dans le vide spatial, elle
  s'enflammerait.
\end{quote}

L'analyse de Lewis et Stalnaker explique cette différence : le monde
le plus proche où je frotte l'allumette dans le vide spatial est
beaucoup plus éloigné du monde actuel que le monde le plus proche où
je me contente de frotter l'allumette. Aussi, bien qu'elle
s'enflamme dans ce dernier monde, elle ne s'enflamme pas dans le
premier. C'est bien le résultat attendu.

\begin{ex}
  La sémantique des sphères invalide cette inférence : on a $p
  \cif r \Entails/ (p \land q) \cif r$. Considérons le modèle
  $\mModel{M} = \tuple{W, O, V}$ où $W = \{w, w_1, w_2\}$,
  $O_w = \{\{w\}, \{w, w_1\}, \{w, w_1, w_2\}\}$,
  $V(p) = \{w_1, w_2\}$, $V(q) = \{w_2\}$ et $V(r) = \{w_1\}$.
  % Précision facultative conservée hors du texte traduit.
%
   Il existe une sphère
  admettant~$p$, à savoir $S = \{w, w_1\}$, et $p \lif r$ est vrai
  dans tous ses mondes ; donc $\mSat{M}{p \cif r}[w]$. Il existe
  aussi une sphère admettant~$(p \land q)$, soit
  $S' = \{w, w_1, w_2\}$, mais
  $\mSat/{M}{(p \land q) \lif r}[w_2]$ ; ainsi,
  $\mSat/{M}{(p \land q) \cif r}[w]$ (voir
  la \olref{fig:antecedent-strengthening}).

\begin{figure}
\begin{center}
\begin{tikzpicture}[layerwidth=1.5,scale=.6]\tiny
  \spheresystem{3}
  \propositionintersect{-10}{3}{30}{5}
  \begin{scope}
    \clip plot[smooth,tension=1] coordinates
          {(0,4.5) (30:.95) (4.5,-3.5)} -- (4.5,4.5) ;
    \spherefill{2}
  \end{scope}
  \draw[proposition] plot[dashed,smooth,tension=1] coordinates
          {(0,4.5) (30:.95) (4.5,-3.5)} ;
  \proposition{30}{2}{30}{5}
  \path (0,0) node[world] {$w$};
  \spherepos{30}{2}{node[world] {$w_1$}}
  \spherepos{-10}{3}{node[world] {$w_2$}}
  \spherepos{-10}{4}{node {$q$}}
  \spherepos{30}{4}{node {$r$}}
  \draw (1,4.5) node[below] {$p$};
\end{tikzpicture}
\caption{Contre-exemple au renforcement de l'antécédent}
\ollabel{fig:antecedent-strengthening}
\end{center}
\end{figure}
\end{ex}

\end{document}
